\section{Recommendations and Outlook}
\label{sec:recommendations}

The results above say what an evaluation has to control for before a content
effect can be claimed, not that earnings calls carry no information. We suggest that future work on this task report six things. A
\textbf{company-disjoint} split next to the temporal one. A \textbf{gap} between
training and test dates of at least the longest horizon. \textbf{Past-volatility
baselines} as the reference for every claim. The \textbf{same-company and random
swaps}, or the released audit command. \textbf{Day~0 measured} from the release
time. And the \textbf{horizon definition} stated with every number. Appendix
Table~\ref{tab:checklist} gives the reason for each and where the release provides
it. \ecvol-bench covers all six for FinCall-Surprise and MAEC, the audit command
applies the checks to any prediction file, and a quickstart script goes from the
released archive of a corpus to the past-volatility baseline, a prediction file
and the audit report in a few seconds, so a new model meets the same floor and
checks on its first day.

There is also reason to expect that content signal can be found once it is
looked for in the right place. Measured release times make
\textbf{event-window targets} possible, such as the volatility of the first
sessions after the release, where language about guidance and surprise should
matter most and where same-company history explains least. \textbf{Structured
features} extracted from the call, such as guidance changes and the gap between
the quantitative surprise and management's framing of it, can be audited one
field at a time, unlike pooled embeddings. They can also be combined with the
fast numeric signal that \citet{fastnumbers2026} show arrives first. \textbf{Audio}
has so far been tested only on lossy webcast transcodes without speaker
diarisation, so recordings at source quality with resolved speakers are the fair
test of prosody. \textbf{Larger company-disjoint corpora} are within reach,
since MM-FinEval \citep{mmfineval2026} is larger than FinCall-Surprise and the
identity recovery used here applies to it. Each of these can be evaluated against the
released baselines on the first day, which is what the protocol is for.
